% Titrage conductimétrique des ions chlorure d'un sérum physiologique (exercice 9, chapitre 3)
\begin{tikzpicture}
\begin{axis}[otMesures,width=104mm,height=56mm,xmin=0,xmax=12.5,ymin=0.1,ymax=0.36,
  xtick={0,2,...,12},ytick={0.1,0.15,...,0.35},xlabel={$V$ (mL)},ylabel={$\sigma$ (mS$\cdot$cm$^{-1}$)},
  every axis x label/.style={at={(axis description cs:1,0)},anchor=north east,yshift=-5mm,font=\small,text=otInk},
  every axis y label/.style={at={(axis description cs:0,1)},anchor=south west,xshift=1mm,font=\small,text=otInk},
  yticklabel style={/pgf/number format/fixed,/pgf/number format/precision=2}]
\solution{\addplot[otGreenDark,line width=0.7pt,domain=0:8] {-0.00187*x+0.17694};
\addplot[otGreenDark,line width=0.7pt,domain=5:12.3] {0.02792*x-0.00546};
\draw[otOrange,line width=0.6pt,dash pattern=on 2pt off 1.5pt] (axis cs:6.12,0.1655) -- (axis cs:6.12,0.1);
\node[pcLegende,anchor=north west,text=otOrange!75!black] at (axis cs:6.3,0.13) {$V_\mathrm{E}\approx6{,}1$ mL};}
\addplot[only marks,mark=*,mark size=1.7pt,otInk] table[x index=0,y index=1] {figures/donnees/03-chlorure-sigma.dat};
\end{axis}
\end{tikzpicture}
